%! Unit 4: Orthogonality — Summative Assessment — Practice Test --- Study Copy (blank)
\documentclass[10pt]{article}
\usepackage{linalg-article}
\usepackage{linalg-boxes}
\newcommand{\bb}{\mathbf{b}}
\newcommand{\xx}{\mathbf{x}}
\newcommand{\pp}{\mathbf{p}}
\newcommand{\ee}{\mathbf{e}}
\newcommand{\zero}{\mathbf{0}}
\newcommand{\avec}{\mathbf{a}}
\newcommand{\ww}{\mathbf{w}}
\newcommand{\uu}{\mathbf{u}}
\newcommand{\T}{^{\mathsf T}}

% Part-divider strip
\newcommand{\parthead}[1]{%
  \vspace{6pt}\noindent
  \begin{headlinebox}{burgundy}{\color{white}\bfseries #1}\end{headlinebox}%
  \vspace{4pt}}

\begin{document}

\pageheader{Unit 4: Orthogonality}{Practice Test --- Study Copy}
\namedateperiod

\vspace{0.10in}
\begin{remindbox}
\small
\textbf{This is a practice test.} It mirrors the real Unit 4 test in length, format,
and the ideas it covers, but the numbers and contexts differ. Work every problem
with no notes, then check your answers against the key.
\end{remindbox}

% ======================================================================
\parthead{Part A --- Vocabulary (8 pts)}
Write the letter of the best description next to each term.

\vspace{4pt}
\noindent
\begin{minipage}[t]{0.40\textwidth}
\begin{enumerate}[leftmargin=2.2em, itemsep=7pt, label=\arabic*.]
  \item Orthonormal vectors \blank{1.2cm}
  \item Orthogonal matrix $Q$ \blank{1.2cm}
  \item Projection matrix $P$ \blank{1.2cm}
  \item Normal equations \blank{1.2cm}
  \item Least squares \blank{1.2cm}
  \item Residual (error) \blank{1.2cm}
  \item Orthogonal complement \blank{1.2cm}
  \item Gram--Schmidt \blank{1.2cm}
\end{enumerate}
\end{minipage}%
\hfill
\begin{minipage}[t]{0.57\textwidth}
\small
\begin{description}[leftmargin=1.4em, itemsep=3pt]
  \item[(A)] The system $A\T A\,\hat\xx=A\T\bb$ whose solution is the best $\hat\xx$.
  \item[(B)] The process that turns independent vectors into an orthonormal set (and builds $A=QR$).
  \item[(C)] Perpendicular \emph{unit} vectors: $q_i\cdot q_j=0$ and each $q_i\cdot q_i=1$.
  \item[(D)] The error vector $\ee=\bb-\pp$; it lies in $N(A\T)$, perpendicular to $C(A)$.
  \item[(E)] A square matrix with orthonormal columns, so $Q\T Q=I$ and $Q\T=Q^{-1}$.
  \item[(F)] All vectors perpendicular to a subspace $V$; together with $V$ it fills the whole space.
  \item[(G)] Choosing $\hat\xx$ to make the error $\|\bb-A\xx\|^2$ as small as possible.
  \item[(H)] The matrix $P=\dfrac{\avec\avec\T}{\avec\T\avec}$ that sends $\bb$ to its projection $\pp=P\bb$.
\end{description}
\end{minipage}

% ======================================================================
\parthead{Part B --- Multiple Choice (12 pts)}
Circle the best answer.

\begin{enumerate}[leftmargin=1.9em, itemsep=9pt]
  \item The least-squares solution makes $A\hat\xx$ the point of $C(A)$ that is
    \hfill (A) farthest from $\bb$ \quad (B) closest to $\bb$ \quad (C) equal to $\bb$ \quad (D) zero
  \item For an orthogonal matrix $Q$, the product $Q\T Q$ equals
    \hfill (A) $0$ \quad (B) $Q$ \quad (C) $I$ \quad (D) $Q^2$
  \item The least-squares error $\ee=\bb-\pp$ is perpendicular to
    \hfill (A) every column of $A$ \quad (B) $\hat\xx$ \quad (C) nothing \quad (D) $\bb$
  \item The normal equations for the best solution of $A\xx=\bb$ are
    \hfill (A) $A\xx=\bb$ \quad (B) $A\T A\,\hat\xx=A\T\bb$ \quad (C) $AA\T\xx=\bb$ \quad (D) $\hat\xx=\bb$
  \item If $Q$ has orthonormal columns, the coordinates of $\bb$ in that basis are found by
    \hfill (A) solving a big system \quad (B) the dot products $q_i\cdot\bb$ \quad (C) inverting $A\T A$ \quad (D) guessing
  \item A projection matrix $P$ always satisfies
    \hfill (A) $P^2=P$ \quad (B) $P^2=I$ \quad (C) $P=0$ \quad (D) $P\T=-P$
\end{enumerate}

% ======================================================================
\parthead{Part C --- Short Answer \& Computation (35 pts)}
Show your work.

\begin{enumerate}[leftmargin=1.9em, itemsep=12pt]
  \item Let $V$ be the line through $\avec=\begin{bsmallmatrix}1\\2\\2\end{bsmallmatrix}$ in $\mathbb{R}^3$.
        Decide whether each vector is in the orthogonal complement $V^{\perp}$ (show the dot product):
        \textbf{(a)} $\ww=\begin{bsmallmatrix}2\\-1\\0\end{bsmallmatrix}$ \quad
        \textbf{(b)} $\uu=\begin{bsmallmatrix}0\\2\\-2\end{bsmallmatrix}$.
        \vspace{2.2cm}
  \item Project $\bb=\begin{bsmallmatrix}3\\1\end{bsmallmatrix}$ onto the line through $\avec=\begin{bsmallmatrix}1\\1\end{bsmallmatrix}$.
        Find $\hat{x}=\dfrac{\avec\cdot\bb}{\avec\cdot\avec}$, the projection $\pp=\hat{x}\avec$, and the error
        $\ee=\bb-\pp$. Verify that $\ee\perp\avec$.
        \vspace{2.4cm}
  \item Use the \textbf{normal equations} $A\T A\,\hat\xx=A\T\bb$ to project
        $\bb=\begin{bsmallmatrix}6\\0\\0\end{bsmallmatrix}$ onto $C(A)$, where
        $A=\begin{bsmallmatrix}1&0\\1&1\\1&2\end{bsmallmatrix}$. Find $\hat\xx$, $\pp=A\hat\xx$, and
        $\ee=\bb-\pp$; check that $A\T\ee=\zero$.
        \vspace{2.8cm}
  \item Fit the best line $y=C+Dt$ to the data $(t,y)=(0,3),(1,3),(2,5),(3,9)$ by least squares.
        Set up $A$ and $\bb$, solve the normal equations for $(C,D)$, and list the residuals
        $\ee=\bb-\pp$. What is $\sum e_i$?
        \vspace{2.8cm}
  \item Let $q_1=\tfrac15\begin{bsmallmatrix}3\\4\end{bsmallmatrix}$ and
        $q_2=\tfrac15\begin{bsmallmatrix}4\\-3\end{bsmallmatrix}$. Show they are \textbf{orthonormal}
        (check $q_1\cdot q_1$, $q_2\cdot q_2$, and $q_1\cdot q_2$), then explain why the matrix
        $Q=\big[\,q_1\ \ q_2\,\big]$ satisfies $Q\T Q=I$.
        \vspace{2.4cm}
  \item Using the same $q_1,q_2$ as above, write $\bb=\begin{bsmallmatrix}5\\0\end{bsmallmatrix}$ in the
        orthonormal basis --- i.e.\ find $c_1,c_2$ with $\bb=c_1q_1+c_2q_2$ --- \emph{without} solving a
        system. (Hint: $c_i=q_i\cdot\bb$.)
        \vspace{2.2cm}
  \item Run \textbf{Gram--Schmidt} on $\avec_1=\begin{bsmallmatrix}3\\4\end{bsmallmatrix}$,
        $\avec_2=\begin{bsmallmatrix}1\\0\end{bsmallmatrix}$ to produce an orthonormal pair $q_1,q_2$.
        (Normalize $\avec_1$; subtract its part from $\avec_2$; normalize the remainder.)
        \vspace{2.6cm}
\end{enumerate}

% ======================================================================
\parthead{Part D --- Extended Response (10 pts)}
Explain your reasoning in full sentences.

\begin{enumerate}[leftmargin=1.9em, itemsep=16pt]
  \item A least-squares line fit produces fitted values $\pp=A\hat\xx$ and residuals $\ee=\bb-\pp$.
        \textbf{(a)} Explain why ``best fit'' forces the error $\ee$ to be perpendicular to the column
        space $C(A)$, so that $A\T\ee=\zero$. \textbf{(b)} When $A$ has an all-ones column, explain why
        this makes $\sum e_i=0$ --- the positive and negative residuals must cancel.
        \vspace{3.0cm}
  \item Let $Q$ be a square matrix with orthonormal columns.
        \textbf{(a)} Explain why the coordinates of any $\bb$ are just the dot products $q_i\cdot\bb$,
        so $\bb=\sum(q_i\cdot\bb)\,q_i$ with no system to solve. \textbf{(b)} Explain why least squares
        collapses to $\hat\xx=Q\T\bb$ when $A=Q$. Use the fact that $Q\T Q=I$.
        \vspace{2.8cm}
\end{enumerate}

\end{document}
